\documentclass[11pt]{amsart}

\usepackage[T1]{fontenc}
\usepackage{lmodern}
\usepackage{microtype}
\usepackage{mathtools,amssymb,amsthm}
\usepackage[hidelinks]{hyperref}

\hypersetup{
  pdftitle={Seven Row-Column Factorial Designs of Strength Two at m=6 Listed as
            Unknown by Rahim and Cavenagh}
}

\newtheorem{theorem}{Theorem}

\emergencystretch=2em

\newcommand{\Idesign}{\mathrm{I}}
\newcommand{\OA}{\mathrm{OA}}
\newcommand{\FF}{\mathbb{F}_2}

\title[Seven row-column factorial designs of strength two]
{Seven Row-Column Factorial Designs of Strength Two at $m=6$ Listed as Unknown
by Rahim and Cavenagh}

\date{}

\subjclass[2020]{05B15, 62K15, 05B20}
\keywords{row-column factorial design, orthogonal array, strength two, Hadamard matrix,
Paley construction, linear code}

\begin{document}

\begin{abstract}
Rahim and Cavenagh conjectured that a Hadamard matrix of order $4m$ forces the existence of an
$\Idesign_k(4m,4n,2,2)$ whenever $n\ge m$, $2^k\mid 16mn$ and $k\le 4m-1$, with the single
exception $m=1$ and $n$ odd. They verified the conjecture for $m\le5$ and recorded that at $m=6$
the unknown cases with minimal parameters are exactly
$\Idesign_{k+3}(24,2^k,2,2)$ for $5\le k\le 11$. We exhibit one witness that settles all seven:
a $24\times 23$ orthogonal array $M$ of strength $2$ from the Paley Hadamard matrix of order $24$,
and a fixed $3$-dimensional binary code $D$ of minimum weight $3$, such that
$A[i][j]=g_i+v_j$ is an $\Idesign_{k+3}(24,2^k,2,2)$, where the $g_i$ are the rows of the first
$k+3$ columns of $M$ and the $v_j$ run over $D^{\perp}$. Both objects are printed in full. The
reading of the source's definition under which these are the objects listed as unknown is the one
fixed by the source's own necessary conditions and by the examples it prints. The same
witness reproduces the authors' already published range $12\le k\le 20$; nothing here settles the
case $m=6$ of the conjecture in full.
\end{abstract}

\maketitle

\section{The conjecture}

Fix integers $q,t\ge2$ and $k\ge t$. An $\OA(N,k,q,t)$ is an $N\times k$ array over $[q]$ in which
every $N\times t$ subarray contains each of the $q^t$ words of $[q]^t$ exactly $N/q^t$ times.
Following Rahim and Cavenagh
\cite[Section~1]{RC2023}, an $\Idesign_k(m,n,q,t)$ is an $m\times n$ array whose $mn$ cells are
filled with the $\lambda$-fold repetition of every word of $[q]^k$, where $\lambda=mn/q^k$, in such
a way that, on fixing any $t$ of the $k$ coordinates, each of the $q^t$ patterns occurs exactly
$n/q^t$ times in every row and exactly $m/q^t$ times in every column. Equivalently, every row is an
$\OA(n,k,q,t)$ and every column is an $\OA(m,k,q,t)$. We only need $q=t=2$, where the condition on a
pair of coordinates $\{a,b\}$ is that each of $00,01,10,11$ occurs $n/4$ times in every row and
$m/4$ times in every column.

The sentence that introduces this definition, in the version of \cite{RC2023} we read and
identically in \cite[p.~79]{Rahim2023thesis}, calls the array formed by a row (column) an orthogonal array ``of
size $k$, degree $n$ (respectively, $m$)'', which interchanges the two slots of the source's own
$\OA(N,k,q,t)$ notation: a row of an $m\times n$ array of words of length $k$ holds $n$ such words,
so as an array it is $n\times k$. Three statements of the source fix the reading, and all three
agree with the definition just given. First, the sentence immediately following it says that fixing
$t$ coordinates inside one row (column) leaves a full factorial design with replication $n/q^t$
(respectively $m/q^t$). Second, the source's own necessary conditions read: if an array of type
$\Idesign_k(m,n,q,t)$ exists then $q^k\mid mn$ and there exist an $\OA(m,k,q,t)$ and an
$\OA(n,k,q,t)$ \cite[p.~81]{Rahim2023thesis}, with $k$ in the degree slot of both. Third, the one
printed example whose shape can settle the orientation is printed with $m$ rows: its
$\Idesign_5(12,8,2,2)$ appears as $12$ rows of $8$ words, the twelve rows marked by the twelve
superscripts $A,\dots,L$ \cite[Table~5.3]{Rahim2023thesis}. The dimensions of the printed
$\Idesign_4(9,9,3,2)$, $\Idesign_4(12,12,2,2)$ and $\Idesign_6(12,16,2,2)$ agree with their labels
as well, though the first two are square and the third is set sideways, so none of them
distinguishes the orientation \cite[Tables~5.1,~5.5,~5.6]{Rahim2023thesis}. Section~3 says which of
this is machine checked.

The statement we settle cases of is Conjecture~5.48 of the first author's thesis
\cite[p.~114]{Rahim2023thesis}, which appears also as a displayed conjecture in \cite{RC2023}. We
quote it from the thesis together with the paragraph that follows it there, which is what makes the
present note finite.

\begin{quote}
\textbf{Conjecture 5.48.} If there exists a Hadamard matrix of order $4m$, then there exists an
$\Idesign_k(4m,4n,2,2)$ for any $n\ge m$, $2^k\mid 16mn$ and $k\le 4m-1$, with the exception $m=1$
and $n$ is odd.

From Theorem~5.39, the above conjecture is true for $m\le5$. For $m=6$, the
unknown cases with minimal parameters are: $\Idesign_{k+3}(24,2^k,2,2)$; $5\le k\le11$. From
Theorem~5.36 and the existence of a Hadamard matrix of order $12$,
$\Idesign_{k+3}(24,2^k,2,2)$ exists for $12\le k\le20$.
\end{quote}

\noindent
Theorems 5.39 and 5.36 are earlier results of the same chapter; the numbering in the published
version of \cite{RC2023} differs.

So the source's list of what is unknown at $m=6$ with minimal parameters has exactly seven
members, and Theorem~\ref{thm:main} below produces all seven at once. Each is an instance of the
conjecture: with $4m=24$ and $4n=2^k$ one has $m=6$, $n=2^{k-2}\ge 8>6$ for $k\ge5$,
$2^{k+3}\mid 16mn=3\cdot 2^{k+3}$, and $k+3\le 14\le 4m-1=23$, while $m=6\ne1$; and a Hadamard
matrix of order $24$ exists, one being written down in Section~\ref{sec:witness}.

\section{The witness}\label{sec:witness}

Let $q=23$ and let $\chi$ be the quadratic-residue character modulo $23$, so $\chi(0)=0$ and
$\chi(a)=\pm1$ according as $a\ne0$ is or is not a square. Index rows and columns of a
$24\times24$ matrix by $0,1,\dots,23$ and set $C_{0,0}=0$, $C_{0,j}=1$ and $C_{i,0}=-1$ for
$i,j\ge1$, and $C_{i,j}=\chi(j-i)$ for $i,j\ge1$. Then $H=C+I$ is a Hadamard matrix of order $24$
\cite{Paley1933}. Negate those rows of $H$ whose entry in column $0$ is $-1$, delete column $0$,
and replace $-1$ by $0$. The result is the $24\times23$ binary array
\begin{center}
\small
\begin{tabular}{r@{\quad}l@{\qquad}r@{\quad}l}
$g_{0}$&\texttt{11111111111111111111111}&$g_{12}$&\texttt{00110101111000001010011}\\
$g_{1}$&\texttt{00000101001100110101111}&$g_{13}$&\texttt{10011010111100000101001}\\
$g_{2}$&\texttt{10000010100110011010111}&$g_{14}$&\texttt{11001101011110000010100}\\
$g_{3}$&\texttt{11000001010011001101011}&$g_{15}$&\texttt{01100110101111000001010}\\
$g_{4}$&\texttt{11100000101001100110101}&$g_{16}$&\texttt{00110011010111100000101}\\
$g_{5}$&\texttt{11110000010100110011010}&$g_{17}$&\texttt{10011001101011110000010}\\
$g_{6}$&\texttt{01111000001010011001101}&$g_{18}$&\texttt{01001100110101111000001}\\
$g_{7}$&\texttt{10111100000101001100110}&$g_{19}$&\texttt{10100110011010111100000}\\
$g_{8}$&\texttt{01011110000010100110011}&$g_{20}$&\texttt{01010011001101011110000}\\
$g_{9}$&\texttt{10101111000001010011001}&$g_{21}$&\texttt{00101001100110101111000}\\
$g_{10}$&\texttt{11010111100000101001100}&$g_{22}$&\texttt{00010100110011010111100}\\
$g_{11}$&\texttt{01101011110000010100110}&$g_{23}$&\texttt{00001010011001101011110}\\
\end{tabular}
\end{center}
which we call $M$. Throughout, a word of $\FF^K$ is a string of $K$ binary digits whose
position $t$, counted from the left starting at $0$, is coordinate $t$; addition is coordinatewise
modulo $2$; and $\langle x,y\rangle$ is the modulo-$2$ inner product. It is a standard consequence
of $HH^{\mathsf T}=24I$ that $M$ is an $\OA(24,23,2,2)$ \cite[p.~148]{HSS1999}: every one of the
$\binom{23}{2}=253$ coordinate pairs sees each of $00,01,10,11$ in exactly $6$ of the $24$ rows.
For $8\le K\le23$ write $G^{(K)}$ for the $24\times K$ array of the first $K$ columns of $M$, whose
rows we call $g^{(K)}_0,\dots,g^{(K)}_{23}$; deleting coordinates preserves strength $2$, so each
$G^{(K)}$ is an $\OA(24,K,2,2)$.

The second ingredient is fixed once and for all. Let
\[
 h_1=\texttt{11100000},\qquad
 h_2=\texttt{01011000},\qquad
 h_3=\texttt{10010100},
\]
regarded inside $\FF^K$ by padding with zeros, and let $D=\langle h_1,h_2,h_3\rangle$. The three are
independent, so $\lvert D\rvert=8$, and the seven nonzero words of $D$ have weights
$3,3,3,3,4,4,4$; in particular
\begin{equation}\label{eq:minwt}
 \text{every nonzero word of } D \text{ has weight at least } 3 .
\end{equation}
Let $V^{(K)}=D^{\perp}\le\FF^K$, a subspace of dimension $K-3$ and hence of size $2^{K-3}$. For
$K=8$ its $32$ words are
\begin{center}
\small
\begin{tabular}{l@{\quad}l@{\quad}l@{\quad}l}
\texttt{00000000}&\texttt{11010000}&\texttt{01101000}&\texttt{10111000}\\
\texttt{10100100}&\texttt{01110100}&\texttt{11001100}&\texttt{00011100}\\
\texttt{00000010}&\texttt{11010010}&\texttt{01101010}&\texttt{10111010}\\
\texttt{10100110}&\texttt{01110110}&\texttt{11001110}&\texttt{00011110}\\
\texttt{00000001}&\texttt{11010001}&\texttt{01101001}&\texttt{10111001}\\
\texttt{10100101}&\texttt{01110101}&\texttt{11001101}&\texttt{00011101}\\
\texttt{00000011}&\texttt{11010011}&\texttt{01101011}&\texttt{10111011}\\
\texttt{10100111}&\texttt{01110111}&\texttt{11001111}&\texttt{00011111}\\
\end{tabular}
\end{center}
(read left to right, top to bottom; the order is immaterial in
Theorem~\ref{thm:main}).
Finally, let $\sigma\colon\FF^K\to\FF^3$ be the syndrome map
$\sigma(x)=(\langle h_1,x\rangle,\langle h_2,x\rangle,\langle h_3,x\rangle)$. Its fibres are the
eight cosets of $V^{(K)}$, and since $h_1,h_2,h_3$ are supported on the first eight coordinates,
$\sigma$ does not depend on $K$. On the $24$ rows of $M$ it takes each of the eight values of
$\FF^3$ exactly three times:
\begin{equation}\label{eq:syn}
 \bigl\lvert\{\,i : \sigma(g_i)=s\,\}\bigr\rvert=3\qquad\text{for every }s\in\FF^3 ,
\end{equation}
equivalently, each of the seven nonzero $h\in D$ satisfies $\langle h,g_i\rangle=1$ for exactly
$12$ of the $24$ rows. Both \eqref{eq:minwt} and \eqref{eq:syn} are finite statements about the
digits printed above.

\begin{theorem}\label{thm:main}
Let $8\le K\le23$, let $g^{(K)}_0,\dots,g^{(K)}_{23}$ be the rows of $G^{(K)}$, let
$v_0,\dots,v_{2^{K-3}-1}$ be the words of $V^{(K)}=D^{\perp}$ in any order, and put
\[
 A[i][j]=g^{(K)}_i+v_j\qquad (0\le i\le23,\ 0\le j\le 2^{K-3}-1).
\]
Then $A$ is an $\Idesign_K\bigl(24,2^{K-3},2,2\bigr)$ of index $3$. Consequently
$\Idesign_{k+3}(24,2^k,2,2)$ exists for every $k$ with $5\le k\le20$; in particular for the seven
values $5\le k\le11$ listed as unknown in the conjecture quoted above.
\end{theorem}

\begin{proof}
Write $n=2^{K-3}$ and $V=V^{(K)}$, $g_i=g^{(K)}_i$, so that $A$ is $24\times n$ and
$24n=3\cdot2^{K}$.

\emph{Multiplicity.} Row $i$ of $A$ lists the coset $g_i+V$, each of its $n$ elements once. For
$x\in\FF^{K}$ we have $x\in g_i+V$ if and only if $\sigma(x)=\sigma(g_i)$, because the fibres of
$\sigma$ are the cosets of $V=\ker\sigma$. Hence $x$ occurs in exactly as many rows as there are
$i$ with $\sigma(g_i)=\sigma(x)$, which is $3$ by \eqref{eq:syn}. So every word of $\FF^K$ occurs
exactly $\lambda=3=24n/2^{K}$ times in $A$.

\emph{Rows.} Fix $i$ and a pair of coordinates $a\ne b$. As $v$ runs over $V$, the entry
$g_i+v$ has $(a,b)$-pattern $\bigl((g_i)_a,(g_i)_b\bigr)+(v_a,v_b)$, so the four pattern counts of
row $i$ are those of the linear map $\pi\colon V\to\FF^2$, $\pi(v)=(v_a,v_b)$, permuted by a
translation of $\FF^2$. Now $\pi$ is surjective unless some nonzero
$\alpha e_a+\beta e_b$ ($\alpha,\beta\in\FF$) annihilates $V$, that is, lies in
$V^{\perp}=D$; and every such word has weight $1$ or $2$, which \eqref{eq:minwt} forbids. So $\pi$
is onto, its fibres are cosets of $\ker\pi$ and each has $\lvert V\rvert/4=n/4$ elements. Every
pattern therefore occurs $n/4$ times in row $i$.

\emph{Columns.} Fix $j$. Column $j$ of $A$ is $\{g_i+v_j : 0\le i\le 23\}$, that is $G^{(K)}$ with
the single constant vector $v_j$ added to every row. Adding a constant permutes the four patterns of
any coordinate pair among themselves, so the pattern counts of column $j$ equal those of
$G^{(K)}$, namely $6=24/4$ each, because $G^{(K)}$ is an $\OA(24,K,2,2)$.

Thus $A$ is an $\Idesign_K(24,n,2,2)$ with $n=2^{K-3}$ and $\lambda=3$. Writing $K=k+3$ turns this
into $\Idesign_{k+3}(24,2^{k},2,2)$, and $8\le K\le23$ is $5\le k\le20$.
\end{proof}

The proof is the specialisation to $\dim D=3$ of a mechanism of \cite{RC2023}, reproduced here so
that the note is self-contained; no originality is claimed for it, nor for the range
$12\le k\le20$, which the quoted passage itself records as following from a published theorem of
\cite{RC2023} together with a Hadamard matrix of order $12$. What was missing was a pair $(M,D)$
satisfying \eqref{eq:minwt} and \eqref{eq:syn} at $4m=24$. The seven cases $5\le k\le11$ are
those the quoted passage lists as unknown, and we claim nothing beyond that listing. Nothing above
proves the case $m=6$ of the conjecture, which asserts an $\Idesign_k(24,4n,2,2)$ for every
admissible pair $(k,n)$ with $n\ge6$; the remaining parameters, in particular
$4n\in\{4b,8b,16b\}$ with $b$ odd, are untouched here, no case $m\ge7$ is settled, and no
nonexistence is proved. We make no claim that the witness above is unique, canonical, or least in
any sense.

On the literature: a search of the arXiv listings for row-column factorial designs and of the
authors' own arXiv postings, made in August 2026, found nothing later than \cite{RC2023} bearing on
these seven cases, and the thesis passage quoted above still states them as unknown. We make no
stronger claim about the literature than that.

\section{What was checked}

Every input of Theorem~\ref{thm:main} is printed in Section~\ref{sec:witness}, so a reader can redo
the argument without a computer. The accompanying program \texttt{verify.py} uses the Python
standard library only and works in exact integer arithmetic. It rebuilds $H$ from $q=23$, checks
all $576$ inner products of its rows and confirms that normalising and reducing $H$ gives the
$24\times23$ array printed above digit for digit; checks that $M$ is an $\OA(24,23,2,2)$; checks
\eqref{eq:minwt} and \eqref{eq:syn}, and that the $32$ printed words of $V^{(8)}$ are exactly
$D^{\perp}$; checks the hypotheses of Theorem~\ref{thm:main} for all sixteen column counts
$8\le K\le23$; and verifies each of the seven designs $5\le k\le11$ cell by cell against the
definition of Section~1, $1\,363\,104$ multiplicity and pattern counts in all. The recorded run
\texttt{verify.output.txt} reports \texttt{VERDICT: ALL 32 CHECKS PASS} and exit status $0$; the
checks it reports beyond those just listed concern statements this note does not make.

Three of the checks concern the reading of the definition. The label and the printed dimensions of
each of the four designs the source exhibits are transcribed into the program, and it checks that
the dimensions match the label, with the rows of the one non-square example the source marks as
rows being the $m$ of that label. It then checks that the reading of Section~1 makes the
replication numbers $n/q^t$, $m/q^t$ and the index $mn/q^k$ positive integers for all four of those
labels and for the seven labels at issue here, while the reading that takes the interchanged slots
literally --- so that the array formed by a row would have $k$ rows and $n$ columns rather than $n$
rows and $k$ columns --- would require $n=k$, which fails for all eleven. Finally it exhibits, for
each of the seven, the two orthogonal arrays that the source's own necessary conditions demand: the
first $K$ columns of $M$ as an $\OA(24,K,2,2)$ and $V^{(K)}$ as an $\OA(2^{K-3},K,2,2)$.

Three limits are worth stating. The cells $12\le k\le20$ are not checked against the definition,
since their arrays have up to $2^{20}$ columns; for those only the hypotheses of
Theorem~\ref{thm:main} are checked, which is what the theorem consumes. The digits of the source's
own printed examples are not transcribed and no run of them is replayed here; only their dimensions
and their parameter arithmetic enter the checks described above. And the passages of the source
quoted or paraphrased in Section~1 were read by us and are not machine checkable at all.

\begin{thebibliography}{9}

\bibitem{HSS1999}
A.~S. Hedayat, N.~J.~A. Sloane, and J.~Stufken,
\emph{Orthogonal Arrays: Theory and Applications},
Springer Series in Statistics, Springer, New York, 1999.

\bibitem{Paley1933}
R.~E.~A.~C. Paley,
\emph{On orthogonal matrices},
J. Math. Phys. \textbf{12} (1933), 311--320.
\href{https://doi.org/10.1002/sapm193312311}{doi:10.1002/sapm193312311}.

\bibitem{Rahim2023thesis}
F.~Rahim,
\emph{Row-column factorial designs and mutually orthogonal frequency rectangles},
Ph.D. thesis, University of Waikato, 2023.
\href{https://hdl.handle.net/10289/15893}{hdl:10289/15893}.
The conjecture quoted in Section~1 is Conjecture~5.48, on p.~114. The definition is on p.~79, the
necessary conditions on p.~81, and the four printed examples in Tables 5.1, 5.3, 5.5 and 5.6, on
pp.~80, 104, 106 and 115.

\bibitem{RC2023}
F.~Rahim and N.~J. Cavenagh,
\emph{Row-column factorial designs with strength at least $2$},
Linear Algebra Appl. (2023),
\href{https://doi.org/10.1016/j.laa.2023.02.018}{doi:10.1016/j.laa.2023.02.018}.
A version of this paper is Chapter~5 of the thesis above. The version of it we read is
\href{https://arxiv.org/abs/2207.02397v3}{arXiv:2207.02397v3}; we quote from the thesis, and the
printed journal wording and numbering may differ from both.

\end{thebibliography}

\end{document}
